\section{Components, scheme structure, and the residual wall}\label{sec:components}
Let $\mathscr D_L$ denote the scheme of maximal minors of $\mu$, and write
\[
 e=Lq,\qquad a=e-p+1,\qquad b=Lc-3,\qquad
 m_X=Lc(k-c),\qquad \mathcal E=\bigoplus_{\nu=1}^L\Sym^2\mathcal S_\nu.
\]
Thus $|\mathscr D_L|=\Delta^{(L)}$. The letters $a,b$ in this section are codimensions, not loading vectors. Theorem~\ref{thm:global-intro} holds throughout its stated range. The structural results below assume the additional stable-range conditions
\begin{equation}\label{eq:stable-range}
 c\ge8,\qquad k\ge2c+1,\qquad e\ge p.
\end{equation}
For the classical input in Lemma~\ref{lem:hankel-strata}, the precise affine dimension statement is $\dim R=2t-2$ on page 2 of \cite{CMSV}'s arXiv version; set $t=r+1$. The secant identification is its equation (2.0.1). Exact-rank strata are open in their rank-at-most closures; confluent forms remain included. These inputs do not supply the component classification below.

These conditions exclude the small equality cases in the rank optimization and the two families of maximal even-dimensional orthogonal Grassmannians.

For distinct $x,y\in\Pj^1$, choose nonzero representatives of their evaluation functionals and let $t\ne0$. Set
\begin{equation}\label{eq:signed-hyperplanes}
 h_+=\operatorname{ev}_x+t\operatorname{ev}_y,\qquad
 h_-=\operatorname{ev}_x-t\operatorname{ev}_y,\qquad
 \ell=\operatorname{ev}_x^{(2n)}-t^2\operatorname{ev}_y^{(2n)}.
\end{equation}
Rescaling the representatives only reparametrizes this family. For
$\varepsilon\in\{+,-\}^L/\{\varepsilon\sim-\varepsilon\}$, let
$\Sigma_\varepsilon$ be the closure of the tuples satisfying
$U_\nu\subset\ker h_{\varepsilon_\nu}$. The all-equal pattern is the common-secant variety $\Sigma_L$ of Proposition~\ref{prop:secant}. Define also
\begin{equation}\label{eq:full-incidence}
 \mathfrak F_L=\{([\ell],(U_\nu)):\rank H_\ell=k,\
                                      H_\ell|_{U_\nu\times U_\nu}=0\},
 \qquad \Phi_L=\overline{\operatorname{pr}_{X_L}\mathfrak F_L}.
\end{equation}
The closure in this definition is essential: the generic annihilator of $\Phi_L$ will be nondegenerate, but its specializations need not be.

\begin{theorem}[Component classification and determinantal scheme]\label{thm:structure}
Assume \eqref{eq:stable-range}.

\smallskip\noindent
\textup{(i)} If $a<b$, the scheme $\mathscr D_L$ is integral and Cohen--Macaulay of codimension $a$, with support $\Phi_L$. Its general point has multiplication corank one, a unique projective annihilator of Hankel rank $k$, and a smooth determinantal neighbourhood. The projection from $\mathfrak F_L$ is birational onto $\Phi_L$.

\smallskip\noindent
\textup{(ii)} If $a=b$, the scheme $\mathscr D_L$ is reduced and Cohen--Macaulay of pure codimension $a$. Its irreducible components are exactly
\[
 \Phi_L\quad\hbox{and}\quad
 \Sigma_\varepsilon\quad
 (\varepsilon\in\{+,-\}^L/\{\pm1\}).
\]
There are $1+2^{L-1}$ components. The general point of each has multiplication corank one and is smooth in the determinantal scheme. The generic annihilator has rank $k$ on $\Phi_L$ and rank two on every $\Sigma_\varepsilon$.

\smallskip\noindent
\textup{(iii)} If $b<a$, the maximal-dimensional irreducible components of $\Delta^{(L)}$ are exactly the $2^{L-1}$ varieties $\Sigma_\varepsilon$, of codimension $b$. Their general points have multiplication corank one and are smooth in $\mathscr D_L$. Every $\Sigma_\varepsilon$ is rational.

\smallskip\noindent
\textup{(iv)} In the expected-codimension range $a\le b$, the fundamental cycle is
\begin{equation}\label{eq:chow-class}
 [\mathscr D_L]=c_a(\mathcal E^*)
 =\left[\prod_{\nu=1}^L\prod_{1\le i\le j\le c}
                         (1+x_{\nu i}+x_{\nu j})\right]_a
       \quad\hbox{in }A^a(X_L),
\end{equation}
where $x_{\nu i}$ are the formal Chern roots of $\mathcal S_\nu^*$ and the brackets select codimension $a$.
\end{theorem}
Part (iii) classifies maximal-dimensional components, not possible smaller components or embedded structure in the excess range. Parts (i) and (ii), by contrast, concern the entire determinantal scheme. Formula~\eqref{eq:chow-class} uses the classical determinantal cycle formula after the expected codimension and reducedness have been established; the cycle formula itself is not a new general intersection-theoretic principle.

\subsection{The two extremal incidences}
\begin{lemma}[Strict rank gap]\label{lem:strict-gap}
Under \eqref{eq:stable-range}, every annihilator incidence stratum has dimension strictly less than $m_X-\min(a,b)$, except possibly the strata
\[
 r=2,\quad s_1=\cdots=s_L=1;
 \qquad
 r=k,\quad s_1=\cdots=s_L=c.
\]
Their dimensions are respectively $m_X-b$ and $m_X-a$.
\end{lemma}
\begin{proof}
Use $D_{r,s}$ and $h_r$ from Lemmas~\ref{lem:hankel-strata} and \ref{lem:isotropic}. If $r=k$, the radical is zero and $s=c$. Suppose $r<k$. For $s=0$, $LD_{r,0}-h_r=(Lc-2)r+1>b$. For $1\le s\le c-2$, the lower endpoint $r=2s$ in the feasible interval gives
\[
 g(s)=1+s(Lc-4)-Ls(s-1)/2.
\]
Its coefficient before minimizing in $r$ is $L(c-s)-2\ge0$. Moreover,
\[
 g(1)=b,\qquad
 g(c-2)-b=(c-3)\{L(c+2)/2-4\}>0.
\]
Concavity places every integer $2\le s\le c-2$ strictly above $b$. At $s=1$ the coefficient $L(c-1)-2$ is strictly positive, so equality requires $r=2$.

For $s=c-1$ and $L=1$, the calculation in Lemma~\ref{lem:optimization} gives at least $a+k-2c+1\ge a+2$. For $s=c-1$ and $L\ge2$, it gives at least
$b+(c-2)\{L(c+1)/2-4\}>b$. For $s=c$, $r<k$, it gives at least $a+1$. Thus, for each fixed nonexceptional $r$ and every feasible $s$, $LD_{r,s}-h_r>\min(a,b)$. Averaging this inequality over the $L$ choices $s_\nu$ gives $\sum_\nu D_{r,s_\nu}-h_r>\min(a,b)$. At $r=2$, replacing any $s_\nu=1$ by zero increases this sum. The two displayed dimensions follow from $D_{2,1}=c$ and $D_{k,c}=q$.
\end{proof}
The proof explains the two ends of the codimension formula: an arbitrary nondegenerate Hankel form imposes $q$ isotropy equations per contact, whereas a rank-two form leaves a choice between two hyperplanes. Intermediate ranks incur a strict dimension loss in \eqref{eq:stable-range}.

\begin{lemma}[Signed secant incidences]\label{lem:signed-secants}
The $2^{L-1}$ varieties $\Sigma_\varepsilon$ are distinct irreducible rational subvarieties of $\Delta^{(L)}$, each of codimension $b$. Their incidence parametrizations are birational. A general point of any one of them admits exactly one rank-two projective annihilator.
\end{lemma}
\begin{proof}
A general point $[h]$ of the secant threefold $S_n$ has a unique unordered secant decomposition. Indeed, evaluations at four distinct points are linearly independent for $n\ge3$, by interpolation; the corresponding confluent statement follows by Hermite interpolation. Thus two different secant lines cannot meet away from their endpoints, and tangent or limiting cases form a proper closed subset. The parameter space is birational to a projective-line bundle over $\operatorname{Sym}^2\Pj^1=\Pj^2$, hence is rational. Sending $[h_+]$ to $[h_-]$ in \eqref{eq:signed-hyperplanes} is a rational involution, independent of the ordering of the two evaluation points.

Normalize the sign pattern by $\varepsilon_1=+$. Over the dense open set of distinct secants with nonzero coefficients, the signed incidence is a product of relative Grassmannians, with dimension $3+Lc(k-1-c)=m_X-b$. It is irreducible and rational: its underlying hyperplane bundles are trivial on a dense open subset of the rational base. The functional $\ell$ in \eqref{eq:signed-hyperplanes} annihilates every product, for either sign, so its image closure lies in $\Delta^{(L)}$.

Here is the projection calculation explicitly. Fix $h\in S_n$ in the above open set. In $\Gr(c,\ker h)$ the exceptional planes admitting another secant hyperplane are the projection of
\[
 \{(h',U):h'\in S_n\setminus\{h\},\
                      U\subset\ker h\cap\ker h'\}.
\]
Its dimension is at most $3+c(k-2-c)$, strictly smaller than $c(k-1-c)$ because $c>3$; if $c>k-2$ it is empty. Therefore a general first plane determines $h$ uniquely. This proves generic injectivity and, in characteristic zero, birationality of the signed projection. The recovered $h$ also determines its partner and the rank-two annihilator. For a general tuple no $U_\nu$ lies in both partner hyperplanes. Consequently two different normalized sign patterns have different images. The exact-rank-two stratum not represented by two distinct evaluations has smaller dimension; it cannot affect these generic assertions.
\end{proof}

\begin{lemma}[The nondegenerate incidence]\label{lem:nondegenerate-incidence}
The variety $\mathfrak F_L$ is smooth and irreducible of dimension $m_X-a$. Its image contains a point whose first plane lies in no secant hyperplane.
\end{lemma}
\begin{proof}
Over the open set of nonsingular Hankel forms, the fibres are products of $\operatorname{OGr}(c,k)$. The restriction equations have independent differentials, as in Lemma~\ref{lem:isotropic}. Since $2c<k$, the special orthogonal group acts transitively on these isotropic Grassmannians; a Witt basis proves transitivity, and the special orthogonal group over $\C$ is connected. Thus the fibres are smooth and geometrically irreducible of dimension $c(k-c)-q$. The base is a smooth irreducible open subset of $\Pj^{p-1}$, so the relative incidence is smooth and irreducible, with the stated dimension.

Fix one nonsingular form $H$ and put $d_H=c(k-c)-q$. In any hyperplane $K$ the restricted form has radical of dimension zero or one. In the nondegenerate case its isotropic $c$-planes have dimension $d_H-c$. In the one-dimensional radical case, the strata avoiding and containing the radical have dimensions respectively
\[
 d_H-c,\qquad d_H-(k-c-1)\le d_H-c,
\]
when nonempty. These formulas follow by the same choice-of-radical-intersection and lift calculation as Lemma~\ref{lem:isotropic}. Varying $K$ among secant hyperplanes contributes only three parameters. Their union in $\operatorname{OGr}(c,k)$ therefore has dimension at most $d_H-c+3<d_H$. Choose $U_1$ outside it and choose the other isotropic planes arbitrarily. This proves the last assertion.
\end{proof}

\subsection{Corank one and the residual multiplication map}
Let $\mathfrak I\subset\Pj(V_{2n}^*)\times X_L$ be the scheme defined by $\ell\mu=0$.
\begin{lemma}[Corank-one chart]\label{lem:incidence-chart}
Above the open subset on which $\mu$ has rank at least $p-1$, the projection from $\mathfrak I$ to $\mathscr D_L$ is an isomorphism. At a point of corank one, infinitesimal pairs satisfy
\begin{equation}\label{eq:incidence-differential}
 \dot\ell(uv)+H_\ell(\dot u,v)+H_\ell(u,\dot v)=0
       \quad(u,v\in U_\nu,\ 1\le\nu\le L).
\end{equation}
Here $\dot U_\nu\in\operatorname{Hom}(U_\nu,V_n/U_\nu)$ and $\dot\ell$ is taken modulo $\C\ell$.
\end{lemma}
\begin{proof}
On a nonzero $(p-1)$-minor, invertible row and column operations reduce $\mu$ to an identity block of size $p-1$ and a remaining row of length $e-p+1$. Its entries generate the maximal-minor ideal. The normalized last coordinate of a projective left kernel vector is one; the other coordinates are uniquely eliminated by the identity block. The remaining equations are exactly that row. This proves the scheme isomorphism, not only a bijection of points. Differentiating the product equations gives \eqref{eq:incidence-differential}.
\end{proof}

\begin{lemma}[Generic smoothness of the signed components]\label{lem:secant-smooth}
If $b\le a$, a general point of each $\Sigma_\varepsilon$ has corank one and the scheme $\mathscr D_L$ is smooth there of codimension $b$.
\end{lemma}
\begin{proof}
The strict rank gap excludes every intermediate-rank annihilator at a general point of $\Sigma_\varepsilon$. If $b<a$, the full-rank incidence has smaller dimension and is excluded as well. If $a=b$ and full-rank annihilators occurred over a dense subset of $\Sigma_\varepsilon$, their fibres would have positive dimension: the projective annihilator space already contains its rank-two point, and its full-rank part is open. This would force $\dim\mathfrak F_L>m_X-a$, a contradiction. Lemma~\ref{lem:signed-secants} then implies corank one. Indeed, any further annihilator would give a projective line of annihilators, and hence either a forbidden rank or more than one rank-two point.

For the tangent calculation take the distinct evaluation points in \eqref{eq:signed-hyperplanes}. Let $g$ be the binary quadratic vanishing at them. Then
\[
 E=\operatorname{rad}H_\ell=gV_{n-2},\qquad
 U_\nu=W_\nu\oplus\C a_\nu,\qquad
 W_\nu=U_\nu\cap E=gK_\nu,
\]
where the $K_\nu$ are independent general $(c-1)$-planes in $V_{n-2}$. The inequality determining this regime is precisely
\begin{equation}\label{eq:residual-wall}
 a-b=L\binom{c}{2}-(2k-5)\ge0.
\end{equation}
Apply Theorem~\ref{thm:global-intro}, already proved in Section~\ref{sec:global-geometry}, with $k'=k-2$ and $c'=c-1$. Its hypotheses hold and \eqref{eq:residual-wall} is its saturation condition. Consequently, on a nonempty open subset of the signed incidence,
\begin{equation}\label{eq:residual-spanning}
 \sum_\nu W_\nu^2=g^2V_{2n-4},\qquad
                  \dim\sum_\nu W_\nu^2=p-4.
\end{equation}
There is no inductive use of the present structural theorem.

Equation~\eqref{eq:incidence-differential} on $W_\nu\times W_\nu$ first requires $\dot\ell$ to annihilate \eqref{eq:residual-spanning}. Thus $\dot\ell$ has dimension $(p-1)-(p-4)=3$. For each such $\dot\ell$, the remaining equations consist of the $c-1$ pairs $(w,a_\nu)$ and the pair $(a_\nu,a_\nu)$. Choose $b_\nu\in V_n$ with $H_\ell(a_\nu,b_\nu)=1$. The $b_\nu$-coordinates of $\dot U_\nu(w)$ and $\dot U_\nu(a_\nu)$ solve these $c$ equations independently; the last coefficient is two, which is nonzero. Thus the tangent dimension of $\mathfrak I$, and by Lemma~\ref{lem:incidence-chart} of $\mathscr D_L$, is
\[
 3+\sum_\nu\{c(k-c)-c\}=m_X-b.
\]
The local dimension is at least this number because the point lies on $\Sigma_\varepsilon$. Equality of dimension and tangent dimension proves regularity of the local ring and the stated smoothness.
\end{proof}
Identity~\eqref{eq:residual-wall} is a deformation-theoretic interpretation of the change of regime. Saturation of the multiplication map after removing the two evaluation points is exactly the condition under which the rank-two component acquires its predicted smooth normal space. This is stronger than locating a minimum in an incidence count.

\subsection{The whole scheme in expected codimension}
\begin{proof}[Proof of Theorem~\ref{thm:structure}]
Suppose first that $a\le b$. On a product of standard affine Grassmannian charts the coordinate ring is a polynomial ring in $m_X$ variables over $\C$, the matrix of $\mu$ has size $p\times e$, and $I=I_p(\mu)$ is its maximal-minor ideal. On every chart meeting the failure scheme, Theorem~\ref{thm:global-intro} gives $\operatorname{height}I\ge a$ and the maximal-minor height inequality gives $\operatorname{height}I\le e-p+1=a$. Since the ring is regular, $\operatorname{grade}I=a$. The Eagon--Northcott grade criterion therefore provides a free resolution of $R/I$ of length $a$; in particular it is Cohen--Macaulay and unmixed of codimension $a$ \cite{EN,Kustin}. This application uses the grade criterion, not a claim that the entries of $\mu$ are independent variables.

Every component consequently has dimension $m_X-a$. By Lemma~\ref{lem:strict-gap}, its generic annihilator must have rank $k$, or, when $a=b$, rank two. In the first case it is dominated by the irreducible incidence $\mathfrak F_L$, so there is at most one such component. Lemma~\ref{lem:nondegenerate-incidence} supplies a point outside every signed secant variety; purity and the strict rank gap force a full-rank component through the resulting non-secant part of the scheme. Hence this component exists and is $\Phi_L$, with dimension $m_X-a$.

The full-rank incidence has this same dimension. Its generic fibre is therefore zero-dimensional, but this fibre is the full-rank open subset of the projective annihilator space. It follows that the generic multiplication corank is one. The generic fibre is a single point, and Lemma~\ref{lem:incidence-chart} identifies a dense open subset of $\Phi_L$ with an open subset of the smooth variety $\mathfrak F_L$. This proves birationality and generic scheme smoothness there.

If $a<b$, the rank-two incidence has smaller dimension, leaving exactly $\Phi_L$. If $a=b$, every distinct $\Sigma_\varepsilon$ has the dimension of a component, and Lemma~\ref{lem:signed-secants} identifies all the rank-two components; the confluent rank-two base is a proper lower-dimensional subset. There are no others by the strict rank gap. Lemma~\ref{lem:secant-smooth} proves generic scheme smoothness on them. A Cohen--Macaulay scheme has no embedded associated primes; if it is generically reduced at every minimal prime, it is reduced. Applying this fact proves (i) and (ii), including integrality in (i).

If $b<a$, all maximal-dimensional components have dimension $m_X-b$. The full-rank incidence and all the intermediate-rank incidences are too small to dominate one. Lemma~\ref{lem:signed-secants} therefore gives exactly the listed components, their codimension and rationality, and Lemma~\ref{lem:secant-smooth} gives their generic corank and smoothness. This proves (iii).

Finally apply the Thom--Porteous formula to the dual map
$V_{2n}^*\otimes\mathcal O\longrightarrow\mathcal E^*$, with source rank $p$, target rank $e$, and rank bound $p-1$ \cite{KL,Fulton}. The expected codimension is $a$, as verified above, so the formula is $c_a(\mathcal E^*)$. The roots of $\Sym^2\mathcal S_\nu^*$ are $x_{\nu i}+x_{\nu j}$, $i\le j$, which gives \eqref{eq:chow-class}. Reducedness shows that in the equality case this is the sum of the component cycles with coefficient one. No excess-range use of the expected-codimension formula is made.
\end{proof}

\subsection{Real native strata and local conditioning}
\begin{corollary}[All signed components have positive native points]\label{cor:signed-native}
For a native interpolation matrix $J$, every $\Sigma_\varepsilon$ has a Zariski-dense set of real points admitting positive, globally separated native realizations. If \eqref{eq:stable-range} holds and $b\le a$, some such points are smooth points of $\mathscr D_L$ of real local codimension $b$. Their pullbacks to loading space have the same local codimension.
\end{corollary}
\begin{proof}
Choose real $x,y>D$ and sufficiently small $t>0$. Every coordinate of $-J\operatorname{ev}_x$ is positive, so both $-Jh_+$ and $-Jh_-$ are strictly positive. Independently for each contact, choose a general $(k-c)$-plane $T_\nu$ containing $-Jh_{\varepsilon_\nu}$ and avoiding the vectors obtained from it by nonconstant coordinate sign changes. The fixed-space construction in Theorem~\ref{thm:fixed-realization} then gives the required globally separated positive loading and $U_\nu\subset\ker h_{\varepsilon_\nu}$. These choices contain a real open subset of the signed incidence. Its birational parametrization is defined over $\R$ and is an isomorphism on dense open subsets. A nonzero real polynomial cannot vanish on this real open subset, so it meets the generic corank-one and residual-spanning opens of Lemma~\ref{lem:secant-smooth}. The real smooth codimension equals the complex codimension there. The loading submersions of Theorem~\ref{thm:fixed-realization} transfer the assertion. In the finite sign test, $s$ and $-s$ define the same unordered coordinate partition. Positivity specifies a Euclidean-open, Zariski-dense real image, not a positivity-defined Zariski-open set.
\end{proof}

\begin{corollary}[Local small-singular-value exponent]\label{cor:local-tail}
At any real smooth corank-one point furnished by Corollary~\ref{cor:signed-native}, there is a loading coordinate neighbourhood, disjoint from the other components, on which the least multiplication singular value $\gamma_{\rm mult}$ satisfies
\begin{equation}\label{eq:local-distance-law}
 C_1\operatorname{dist}(A,|\mathscr D_L|_{\rm load})
 \le\gamma_{\rm mult}(A)\le
 C_2\operatorname{dist}(A,|\mathscr D_L|_{\rm load}).
\end{equation}
Here the neighbourhood has compact closure in a fixed frame chart and in a fixed full-rank loading sign chamber. On a sufficiently small coordinate product box, a probability law with density bounded above and below by positive constants satisfies
\begin{equation}\label{eq:local-tube-tail}
 c_1\epsilon^b\le\Pr\{\gamma_{\rm mult}\le\epsilon\}\le c_2\epsilon^b
       \quad(0<\epsilon<\epsilon_0).
\end{equation}
At a real smooth corank-one point of $\Phi_L$ in the range $a\le b$, the same assertion holds with $a$ in place of $b$.
\end{corollary}
\begin{proof}
Use the Schur-complement row of Lemma~\ref{lem:incidence-chart}. Bounded invertible matrix operations compare its Euclidean norm with $\gamma_{\rm mult}$. At a smooth scheme point of codimension $b$, the differentials of its entries span a conormal space of dimension $b$. Select $b$ independent entries as normal coordinates. The whole row vanishes on the failure submanifold; bounded derivatives give the upper distance estimate, and the selected coordinates give the lower estimate. The loading submersion and compact chart bounds preserve these comparisons. Integration in the $b$ normal coordinates, with bounded Jacobian and density on the product box, gives \eqref{eq:local-tube-tail}. The full-rank case is identical with codimension $a$.
\end{proof}
These are local probability laws near specified smooth strata. They do not assert a global tail for arbitrary random loadings or constants uniform in $c,k,L$. In particular, singular intersections and smaller components in the excess range are not covered by this local tube calculation.
